\section{Limitations}\label{sec:limitations}

\paragraph{The start is assumed, not measured.} Every empirical statement in
this paper is conditional on the dive credit $c$. The credit is anchored to
elite 15~m start times, and the band we sweep was widened twice, by logged
amendment, once we converted those times to this field's slower pace; but
no source in our data publishes a 15~m split, so the credit is never
observed for any race we analyse. Because the credit and both positive-split
couplings push lap~1 the same way, the fitted $\beta_x$ and $\gamma$ are
credit-conditional estimates, and the width of the credit band is the width
of our ignorance about them. A pace-dependent start term,
$c(v) \approx 15/v - t_{15}$, and 15~m timing at even one meet would do more
for the model comparison than any additional race sheet.

\paragraph{Turns and underwaters are folded into $\phi$.} Short-course yards
racing contains seven turns in 200 yards, and the underwater phase is faster
than surface swimming. The model represents the whole course discount by one
calibrated factor $\phi$ and treats each 50 as free swimming at constant
velocity. The shape-invariance theorem means $\phi$ cannot move the M3
optimum, so this does not bias the shape comparison, but it does mean the
model cannot distinguish a swimmer who fades because of surface economy from
one whose turns deteriorate; both appear as $\beta_x$.

\paragraph{Constant velocity within a 50.} Cost is convex in velocity, so
this systematically understates energy cost, and the missing cost is absorbed
into the calibrated engine. It also excludes by assumption the within-lap
oscillation mechanisms of the optimal-control literature; split-level data
cannot test those in either direction.

\paragraph{One region, one course, a broad field.} All eleven meets are in
Northern California; the sample is one competitive culture and a handful of
pools. Final times span 1:37 to 2:40, a developmental range in which pacing
skill itself varies, and the population rule admits 54 races on club-linked
rather than published age. We report the published-age stratum separately
and the conclusions do not change, but a replication elsewhere, on
long-course metres, and on female and older swimmers is required before any
of this is a statement about the 200 freestyle in general.

\paragraph{Preliminary swims and the finishing kick.} Fifty-seven usable
races are preliminaries, which may be paced to qualify rather than to
minimize time. And the observed mean profile has a feature none of the five
models can produce: the third lap is the slowest and the fourth is faster
than it. Every model in the family predicts monotone lap times, because
every mechanism is monotone in position or reserve. A terminal kick is the
signature of a swimmer holding something back and spending it once the
finish is certain, which is a rational response to uncertainty about one's
own state or a tactical response to the field, and neither is in the model.

\paragraph{Fitted parameters are not physiology.} $\beta_x$, $\gamma$ and
$\beta_E$ are estimates under this model family, conditional on the start
credit, on the constant-velocity assumption and on the calibrated engine.
They are not measurements of any swimmer, and we have taken care nowhere to
describe them as such. The same applies with more force to the engine
parameters $E_0$, $R$ and $\phi$, which are calibrated to reproduce a target
time and, by Theorem~\ref{thm:shape}, are not identifiable from pacing data
at all.

\paragraph{Expectation is weak.} The pre-race personal best is derived only
from earlier rows in this dataset, so for most swimmers it is one earlier
race, often months before, in a developing athlete. Improvement on that
number mixes training gains with pacing quality, which is one reason the
deviation--performance regression was pre-registered as likely underpowered
and why its result is reported as a null rather than explained away.

\paragraph{What we did not do.} We did not fit any engine or hull parameter
to pacing data, did not use AIC or likelihood comparison, did not examine
the test set before the registered evaluation, and did not adjust anything
after it. The exploratory 80-race pilot informed the start-credit band and
exposed the sign error; both are logged as amendments made before any
calibration. We regard the ordering of those steps as the main
methodological contribution of the empirical half of the paper.
